\documentclass[12pt,twoside]{article}
%\usepackage[table]{xcolor} % important to avoid options clash.
%\input{02465shared_preamble}
%\usepackage{cleveref}
\usepackage{url}
\usepackage{graphics}
\usepackage{multicol}
\usepackage{rotate}
\usepackage{rotating}
\usepackage{booktabs}
\usepackage{hyperref}
\usepackage{pifont}
\usepackage{latexsym}
\usepackage[english]{babel}
\usepackage{epstopdf}
\usepackage{etoolbox}
\usepackage{amsmath}
\usepackage{amssymb}
\usepackage{multirow,epstopdf}
\usepackage{fancyhdr}
\usepackage{booktabs}
\usepackage{xcolor}
\newcommand\redt[1]{ {\textcolor[rgb]{0.60, 0.00, 0.00}{\textbf{ #1} } } }


\newcommand{\m}[1]{\boldsymbol{ #1}}
\newcommand{\EE}{\mathbb{E}}
\DeclareMathOperator*{\argmax}{arg\,max}
\DeclareMathOperator*{\argmin}{arg\,min}
\newcommand{\yoursolution}{ \redt{(your solution here) } } 



\title{ Report 4 hand-in }
\date{ \today }
\author{Alice (\texttt{s000001})\and  Bob (\texttt{s000002})\and Clara (\texttt{s000003}) } 

\begin{document}
\maketitle

\begin{table}[ht!]
\caption{Attribution table. Feel free to add/remove rows and columns}
\begin{tabular}{llll}
\toprule
                                    & Fabian   & Valdemar   \\
\midrule
 1: UCB-based exploration           & 60\%  & 40\%  \\
 2: Linear function approximator    & 40\%  & 60\%  \\
 3: Quadratic function approximator & 40\%  & 60\%  \\
\bottomrule
\end{tabular}
\end{table}

%\paragraph{Statement about collaboration:}
%Please edit this section to reflect how you have used external resources. The following statement will in most cases suffice: 
%\emph{The code in the irls/project1 directory is entirely}

%\paragraph{Main report:}
Headings have been inserted in the document for readability. You only have to edit the part which says \yoursolution. 

\section{Finding the rebels using UCB-exploration (\texttt{rebels.py})} 

\section{Computing State-action values using random probe-droids (\texttt{probe\_droids.py})}

\subsubsection*{{\color{red}Problem 2: Linear function approximator}}

We implement a linear feature encoder that maps gridworld states to feature vectors. Given a state $s = (i, j)$ where $i$ and $j$ are the row and column, we compute normalized coordinates:
\begin{align}
u &= \frac{i}{H-1}, \quad v = \frac{j}{W-1}
\end{align}
where $H$ is the grid height and $W$ is the grid width.

The linear feature vector is:
\begin{align}
z_{\text{lin}}(s) = \begin{bmatrix} 1 \\ u \\ v \end{bmatrix}
\end{align}

For the action-dependent feature representation, we construct a $4\ell$-dimensional vector where $\ell = 3$ is the dimension of $z_{\text{lin}}(s)$. Each action $a \in \{0, 1, 2, 3\}$ gets its own 3-dimensional block:
\begin{align}
x(s, a) = \begin{bmatrix} \mathbf{0}_{a} \\ z_{\text{lin}}(s) \\ \mathbf{0}_{a'} \end{bmatrix}
\end{align}

The Q-values are then approximated as $Q(s,a) = x(s,a)^\top w$, where $w \in \mathbb{R}^{12}$ is the weight vector.

We train a LinearSemiGradSarsa agent with this feature encoder using a uniform random policy ($\epsilon = 1$) for 20,000 episodes with learning rate $\alpha = 0.02$ and discount factor $\gamma = 1.0$.

\subsubsection*{{\color{red}Problem 3: Quadratic function approximator}}

For the quadratic approximator, we augment the linear features with quadratic terms and interaction terms:
\begin{align}
z_{\text{quad}}(s) = \begin{bmatrix} 1 \\ u \\ v \\ u^2 \\ v^2 \\ uv \end{bmatrix}
\end{align}

This gives a state feature dimension of $\ell = 6$. The action-dependent feature vector $x(s,a)$ similarly has dimension $d = 4\ell = 24$, with the quadratic feature vector placed in the block corresponding to action $a$.

We train a LinearSemiGradSarsa agent with this feature encoder using the same settings as the linear case: uniform random policy, 20,000 episodes, $\alpha = 0.02$, and $\gamma = 1.0$.

The quadratic feature approximator can capture more complex value functions than the linear approximator, potentially providing better Q-value estimates despite using the same amount of training data.

\end{document}